\chapter{Additional Code Listings}\label{appendix:code-listings}

This appendix provides the lower-level rendering and SVG assembly workflow summarized in Chapter~\ref{chap:DesignAndImplementation}. The public composition script calls these operations through \texttt{Visualizer.compose()}.

\begin{lstlisting}[caption={Internal SVG composition workflow},label={lst:svg-layers-example}]
# Internal workflow (called by compose())

# For each panel: render layers, export to SVG
panel_viz = Visualizer(audio=audio, score=score,
                       maps=maps, beats=beats)
panel_viz.add_panel(MelSpec(), BeatLogits())
panel_viz.load_all_layers(audio_path=audio,
                          beat_file=beats)
fig, ax = panel_viz.draw()
panel_svg = panel_viz.turn_to_SVG(
  filename="panel.svg",
  svg_warped_score=score)

# Stack panels with score using create_final_SVG()
composer = Visualizer(audio=audio, score=score,
                      maps=maps, beats=beats)
composer.create_final_SVG(
  svg_layers=[(panel_svg, 50),
              (score_svg, 300)],
  output_file="final.svg")
\end{lstlisting}

Each panel is rendered to SVG through \texttt{turn\_to\_SVG()}, preserving its layers as named SVG groups. The resulting panel SVGs are then stacked with the warped score by \texttt{create\_final\_SVG()}, with vertical offsets controlling their positions while the shared timeline preserves horizontal alignment.